% Author: SuFan; MIT License. All values are illustrative.
\documentclass{article}
\usepackage[margin=20mm]{geometry}
\usepackage{gradbars,booktabs}
\gradbarsrequireversion{1.1.1}
\pagestyle{empty}
\begin{document}
% BEGIN TABLE
\begingroup
\gradbarscategory{Desk}{blue}{diagonal}
\gradbarscategory{Locker}{teal}{dots}
\gradbarscategory{Delivery}{orange}{crosshatch}
\gradbarssetup{print,min=0,max=300,width=70mm,height=14pt,
  stack names={Desk,Locker,Delivery},segment labels=percent,
  precision=0,label=none,table legend=true}
\renewcommand{\arraystretch}{1.8}
\begin{center}
\begin{tabular}{@{}lrl@{}}
\toprule
Branch & Loans & \gradtablelegendhead[print]
  {Fulfilment share}{Desk,Locker,Delivery}\\
\midrule
North & 300 & \gradstack{180,90,30}\\
South & 300 & \gradstack{120,90,90}\\
East & 300 & \gradstack{60,180,60}\\
West & 300 & \gradstack{150,105,45}\\
\bottomrule
\end{tabular}
\end{center}
\endgroup
% END TABLE
\noindent Each row contains 300 loans. Inputs are counts; segment labels divide
by that row's total. Equal full widths follow from these equal totals, not from
automatic normalization. Textures keep the three channels distinct in print.
\end{document}
